\begin{lemma}[Extension-base Part I(1) and the support-class order]
Let \(T\) be an \(n\)-rung ETT, let its terminal iteration \(n\) be a PE
with supporting vertex \(v_n\), and put \(j=m(v_n)\).  Let \(L\) be the
active index set from \noderef{ExtensionBaseActiveIndices}.  Assume the
following explicit lower-rung instance of technical assertion (v).  For
every \(a<r\leq n\), every proper coloring \(\pi\) leaving exactly \(e_0\)
uncolored and \((T_r,D_{r-1},\varphi_{r-1})\)-stable, every
defective color \(\delta\) of \(T_r\) under \(\pi\), and every
non-first vertex \(x\) of
(I[\partial_{\pi,\delta}(T_r)]), if (m(x)\leq a), then

\[
 x\preceq v_a
\]

for each PE or SE supporting/extension vertex (v_a).  In Lean,
``non-first'' is supplied by an explicit earlier boundary-color vertex,
and (x\preceq v_a) means equality or
\noderef{TreeVertexOrder} in (T_r).

Then the following statements hold.

\begin{enumerate}
\item If (s,t\in L) and (s<t), then (v_t\preceq v_s) in the
      displayed order of (T).
\item For every (s\in L),
      (m(v_s)=j) and (v_n\preceq v_s).
\item If (s,t\in L), (v_s\ne v_t), and (v_s\prec v_t), then
      (t<s).  Consequently, after the distinct support vertices are
      named (w_1\prec\cdots\prec w_\kappa), every index in the class
      of (w_h) is larger than every index in the class of (w_i)
      whenever (h<i), and (v_n=w_1).
\end{enumerate}
These are Part I(1) and precisely the class-order consequences used from
Part I(3).
\end{lemma}
\begin{proof}
Fix (s,t\in L) with (s<t).  Membership in
\noderef{ExtensionBaseActiveIndices} gives
\[
 j\leq s<t,\qquad \Theta_s=\Theta_t=\mathrm{PE},qquad
 m(v_t)\leq j.                                                   \tag{1}
\]
Inspect the certificate at iteration (t), supplied by
\noderef{TashkinovSeries}.  The equality
\(\Theta_t=\mathrm{PE}\) eliminates the initial, SE, and RE constructors
of \noderef{TashkinovSeriesStep}.  In its PE constructor, write
\(\pi_{t-1}\) for the selected coloring, (f_t) for the selected edge,
and \(\delta_t\) for its color.  The fields of
\noderef{MaximumDefectiveChoice} say that \(\pi_{t-1}\) is proper,
leaves precisely (e_0) uncolored, is
((T_t,D_{t-1},\varphi_{t-1}))-stable, and that (f_t) is a defective
boundary edge incident with (v_t), colored \(\delta_t\).

Expanding \noderef{DefectiveEdge} and \noderef{DefectiveColor}, the
\(\delta_t\)-boundary contains at least two distinct edge identities.
Choose one different from (f_t), and let (x) be its endpoint in
(T_t).  Properness, in the sense of \noderef{IsProperEdgeColoring},
forces (x\ne v_t): otherwise the two distinct incident edges would
have the same color.  Hence (x) and (v_t) are distinct
\(\delta_t\)-boundary vertices, as defined by
\noderef{BoundaryColorVertex}.  Each is a defective vertex by
\noderef{DefectiveVertex}.  The maximality field of
\noderef{MaximumDefectiveChoice}, applied to (x) and the same selected
coloring, gives (x\prec_{T_t}v_t).  Thus (v_t) is not the first
vertex of (I[\partial_{\pi_{t-1},\delta_t}(T_t)]).

By (1), (m(v_t)\leq j\leq s).  Apply the stated lower-rung assertion
(v) to the ETT ending at (T_t), the coloring \(\pi_{t-1}\), the
defective color \(\delta_t\), the non-first vertex (v_t), and the PE
support (v_s).  It yields

\[
                         v_t\preceq_{T_t}v_s.                    \tag{2}
\]

Every certified step makes its old tree an ordered initial segment of
its new tree by \noderef{TashkinovSeriesStep} and
\noderef{TreeSequenceInitialSegment}; the two terminal segment fields of
\noderef{ExtendedTashkinovTree} then embed (T_t) order-preservingly in
(T).  Therefore (2) is also (v_t\preceq_Tv_s).  This proves the
first assertion.

The terminal PE support (v_n) lies in (T_n) by the defective-vertex
field of its PE certificate.  Hence (j=m(v_n)\leq n), and the terminal
index belongs to (L).  For (s\in L), either (s=n), when the desired
order is equality, or the first assertion with (t=n) gives
\(v_n\preceq_Tv_s\).  The definition
\noderef{ETTFirstAppearanceIndex} takes the least ladder stage containing
a vertex.  All ladder stages are ordered initial segments of (T), as
just observed.  Thus a stage containing (v_s) also contains every
vertex preceding (v_s), in particular (v_n); consequently
\(m(v_n)\leq m(v_s)\).  Active membership gives the reverse inequality
\(m(v_s)\leq j=m(v_n)\).  Hence (m(v_s)=j), proving the second
assertion.

Finally take distinct active supports (v_s,v_t) with
\(v_s\prec_Tv_t\).  If (s=t), the exact chosen-vertex field at that
single PE slot gives (v_s=v_t), contrary to the hypothesis.  If
\(s<t\), the first assertion gives \(v_t=v_s\) or
\(v_t\prec_Tv_s\); both alternatives contradict respectively
distinctness or asymmetry of the strict order induced by the
duplicate-free vertex list in \noderef{TreeVertexOrder}.  Trichotomy of
the natural indices therefore leaves only (t<s).

Order the distinct supports as
\(w_1\prec\cdots\prec w_\kappa\).  Applying the preceding conclusion to
one index from the classes of (w_h,w_i), with (h<i), shows that every
index of the earlier vertex (w_h) is larger than every index of
(w_i).  Since the terminal index is larger than every other active
index, its support must be the earliest class vertex (w_1).  This is
exactly the asserted consequence of Part I(3).
\end{proof}
